\booksection{Lampiran Bab 11: Pengetahuan Teknik Imaging Profesional}

\chapterintro{Untuk benar-benar memahami misteri fotografi dan video ponsel, Anda perlu memahami beberapa pengetahuan teknik imaging profesional. Pengetahuan ini tidak hanya membantu Anda menggunakan kamera ponsel dengan lebih baik, tetapi juga membuat Anda lebih bijak saat memilih perangkat.}

\booksubsection{Resolusi, Frame Rate, dan Bitrate: Tiga Elemen Kualitas Video}

\subsubsection*{Resolusi menentukan kejernihan gambar, dinyatakan dalam jumlah piksel.}

720P (HD)

\begin{itemize}
  \item \textbf{: 1280×720 piksel, kejernihan tingkat pemula.}1080P (Full HD)
  \item \textbf{: 1920×1080 piksel, resolusi paling mainstream saat ini.}2K (QHD)
  \item \textbf{: 2560×1440 piksel, lebih jernih dari 1080P.}4K (UHD)
  \item \textbf{: 3840×2160 piksel, ultra HD, detail kaya.}8K
  \item \textbf{: 7680×4320 piksel, standar masa depan, saat ini dukungan perangkat terbatas.}Semakin tinggi resolusi, semakin jernih gambar, tetapi file juga semakin besar, dan persyaratan penyimpanan serta pemrosesan semakin tinggi.
\end{itemize}

Frame Rate:

\subsubsection*{Frame rate menentukan kelancaran gambar, dinyatakan dalam frame per detik (fps).}

24fps

\begin{itemize}
  \item \textbf{: Frame rate standar film, memiliki "cinematic feel" yang unik.}30fps
  \item \textbf{: Umum digunakan untuk TV dan video online, lancar dan natural.}60fps
  \item \textbf{: Frame rate tinggi, cocok untuk scene olahraga dan recording game, sangat lancar.}120fps/240fps
  \item \textbf{: Frame rate super tinggi, dapat digunakan untuk video slow motion.}Semakin tinggi frame rate, semakin lancar gambar, semakin baik efek slow motion, tetapi file juga semakin besar.
\end{itemize}

Bitrate:

\subsubsection*{Bitrate menentukan kepadatan data video, dinyatakan dalam data per detik (Mbps).}

Semakin tinggi bitrate, semakin baik kualitas gambar, semakin banyak detail yang dipertahankan, tetapi file juga semakin besar.

\begin{itemize}
  \item Bitrate terlalu rendah akan menyebabkan gambar muncul mozaik, color banding, dan jejak kompresi lainnya.
  \item Shooting video profesional biasanya memerlukan bitrate tinggi untuk mempertahankan lebih banyak ruang post-processing.
  \item 2. Format Kodek Video: H.264, H.265, dan ProRes
\end{itemize}

\booksubsection{Format Kodek Video: H.264, H.265, dan ProRes}

\subsubsection*{Karakteristik}

\begin{itemize}
  \item \textbf{Karakteristik}Kelebihan
  \item \textbf{: Ruang penyesuaian post-processing sangat besar, dapat mengubah eksposur, white balance, dan parameter lain tanpa kehilangan kualitas gambar.}Kekurangan
  \item \textbf{Kekurangan}Skenario Cocok
  \item \textbf{: Produksi profesional kelas atas. Seri iPhone 15 Pro mendukung shooting video Apple ProRes (hingga 4K 60fps), tetapi format ProRes RAW saat ini perlu digunakan bersama perangkat perekam eksternal, ponsel belum mendukung shooting langsung.}H.265 (HEVC):
\end{itemize}

\subsubsection*{Karakteristik}

\begin{itemize}
  \item \textbf{Karakteristik}Kelebihan
  \item \textbf{: Ruang penyesuaian post-processing sangat besar, dapat mengubah eksposur, white balance, dan parameter lain tanpa kehilangan kualitas gambar.}Kekurangan
  \item \textbf{Kekurangan}Skenario Cocok
  \item \textbf{: Produksi profesional kelas atas. Seri iPhone 15 Pro mendukung shooting video Apple ProRes (hingga 4K 60fps), tetapi format ProRes RAW saat ini perlu digunakan bersama perangkat perekam eksternal, ponsel belum mendukung shooting langsung.}Apple ProRes:
\end{itemize}

\subsubsection*{Karakteristik}

\begin{itemize}
  \item \textbf{Karakteristik}Kelebihan
  \item \textbf{: Ruang penyesuaian post-processing sangat besar, dapat mengubah eksposur, white balance, dan parameter lain tanpa kehilangan kualitas gambar.}Kekurangan
  \item \textbf{Kekurangan}Skenario Cocok
  \item \textbf{: Produksi profesional kelas atas. Seri iPhone 15 Pro mendukung shooting video Apple ProRes (hingga 4K 60fps), tetapi format ProRes RAW saat ini perlu digunakan bersama perangkat perekam eksternal, ponsel belum mendukung shooting langsung.}ProRes RAW:
\end{itemize}

\subsubsection*{Karakteristik}

\begin{itemize}
  \item \textbf{Karakteristik}Apa itu RAW
  \item \textbf{: Data sensor mentah yang belum diproses, mempertahankan informasi maksimal.}Kelebihan
  \item \textbf{: Ruang penyesuaian post-processing sangat besar, dapat mengubah eksposur, white balance, dan parameter lain tanpa kehilangan kualitas gambar.}Kekurangan
  \item \textbf{Kekurangan}Skenario Cocok
  \item \textbf{: Produksi profesional kelas atas. Seri iPhone 15 Pro mendukung shooting video Apple ProRes (hingga 4K 60fps), tetapi format ProRes RAW saat ini perlu digunakan bersama perangkat perekam eksternal, ponsel belum mendukung shooting langsung.}: Produksi profesional high-end. Seri iPhone 15 Pro dan setelahnya mendukung syuting video Apple ProRes (maksimal 4K 60fps), iPhone 17 Pro lebih lanjut mendukung ProRes RAW langsung rekam di kamera.
\end{itemize}

\booksubsection{HDR: Teknologi High Dynamic Range}

HDR10

\begin{itemize}
  \item \textbf{: Standar HDR paling dasar, kompatibilitas baik.}HDR10+
  \item \textbf{: Standar yang didorong Samsung, mendukung metadata dinamis.}Dolby Vision
  \item \textbf{: Standar HDR kelas atas yang dikembangkan Dolby, kualitas gambar terbaik, tetapi memerlukan lisensi.}HLG
  \item \textbf{: Standar HDR tingkat broadcast, kompatibel dengan monitor biasa.}Saat shooting video HDR, diperlukan monitor yang mendukung HDR untuk melihat efek lengkap.
\end{itemize}

4. Color Space dan Color Depth

\booksubsection{Color Space dan Color Depth}

\subsubsection*{sRGB}

\begin{itemize}
  \item \textbf{: Color space standar, kompatibilitas terbaik.}Display P3
  \item \textbf{: Standar wide gamut yang dipromosikan Apple, warna lebih kaya.}Rec.2020
  \item \textbf{: Color space standar video HDR, gamut terluas.}Color Depth:
\end{itemize}

\subsubsection*{8-bit}

\begin{itemize}
  \item \textbf{: 16,7 juta warna, standar video biasa.}10-bit
  \item \textbf{: 1,07 miliar warna, standar video profesional, transisi warna lebih halus.}12-bit
  \item \textbf{: 68 miliar warna, standar profesional kelas atas.}Saran Praktis
\end{itemize}

\begin{tcbinfo}{Shooting harian: 1080P/4K + 30fps + encoding H.265, menyeimbangkan kualitas gambar dan ukuran file.}
\begin{itemize}
  \item Shooting olahraga: 1080P/4K + 60fps, memastikan kelancaran.
  \item Kreasi profesional: 4K + ProRes (jika didukung), mempertahankan ruang post-processing maksimal.
  \item Saran penyimpanan: Video bitrate tinggi sangat menghabiskan ruang, disarankan cadangkan secara berkala ke komputer atau cloud storage.
  \item Jelajahi: Pengaturan Video Ponsel Anda
\end{itemize}

\end{tcbinfo}

\begin{tcbpractice}{Buka aplikasi kamera, beralih ke mode video.}
\begin{enumerate}
  \item Lihat opsi resolusi dan frame rate, coba berbagai kombinasi.
  \item Jika didukung, periksa apakah ada opsi lanjutan seperti HDR, ProRes.
  \item Rekam beberapa segmen video dengan pengaturan berbeda, bandingkan kualitas gambar dan ukuran file.
  \item Bab 12: Hiburan Audio-Visual dan Membaca
\end{enumerate}

\end{tcbpractice}

\booksubsection{Perbandingan Format RAW dengan HEIF/JPEG}

\subsubsection*{Apa itu format RAW:}

Format RAW adalah “data mentah” yang direkam sensor kamera — informasi mentah murni tanpa kompresi dan pemrosesan. Ibarat bahan mentah, perlu “dimasak” di post-production agar jadi hasil akhir. Karena penyimpanan dan daya komputasi ponsel terbatas, default foto yang diambil sudah “dimasak” menjadi JPEG atau HEIF, RAW perlu diaktifkan khusus.

\subsubsection*{Perbandingan RAW vs HEIF vs JPEG:}

\begin{itemize}
  \item \textbf{JPEG}: Format tertua, kompatibilitas terbaik (semua perangkat bisa membuka). Color depth 8-bit, kompresi lossy, file terkecil (2-5MB per foto). Ruang post-production terbatas, paksa menerangi bayangan mudah muncul color banding.
  \item \textbf{HEIF/HEIC}: Format baru yang didorong Apple, dengan kualitas yang sama file hanya setengah dari JPEG (1-3MB per foto). Mendukung color depth 10-bit dan channel transparan. Kelemahannya beberapa perangkat lama dan Windows perlu install plugin untuk membuka.
  \item \textbf{RAW}: Data mentah tanpa kompresi, file terbesar (20-100MB per foto). Color depth biasanya 12-16bit, ruang post-production terbesar. Kelemahannya harus pakai software profesional (Lightroom, Snapseed, dll.) untuk membuka dan memproses, tidak bisa langsung dibagikan.
\end{itemize}

\subsubsection*{Mengapa fotografer profesional memotret RAW:}

\begin{itemize}
  \item \textbf{White balance bisa disesuaikan tanpa kerugian di post-production}: Jika white balance JPEG salah (misalnya lampu pijar kuning), koreksi post-production akan mengurangi kualitas; RAW bisa reset suhu warna, sepenuhnya tanpa kerugian.
  \item \textbf{Toleransi eksposur besar}: Bayangan JPEG underexpose 2 stop pada dasarnya hitam pekat, setelah diterangkan penuh noise; RAW bisa "menyelamatkan" detail usable dari bayangan underexpose 3-4 stop. Overexpose sama.
  \item \textbf{Detail bayangan dan highlight lebih kaya}: Dynamic range RAW jauh lebih besar dari JPEG, ruang highlight tidak overexpose dan bayangan tidak dead black jauh lebih besar.
  \item \textbf{Sharpening, noise reduction, koreksi lensa semua terkontrol di post-production}: Pemrosesan in-kamera adalah "sekali selesai", RAW bisa diatur berulang, setiap kali mulai dari data mentah.
\end{itemize}

\subsubsection*{Cara mengaktifkan RAW di ponsel:}

\begin{itemize}
  \item \textbf{iPhone ProRAW}: Hanya iPhone 12 Pro dan model Pro di atasnya yang mendukung. Buka "Settings Camera Format" aktifkan "Apple ProRAW". Lalu di antarmuka kamera atas akan muncul tombol "RAW", aktifkan saja. File sekitar 25MB/foto, disarankan aktifkan saat dibutuhkan.
  \item \textbf{Android Motion Raw / Ultra RAW}: Setiap merek penyebutannya berbeda. Model flagship Xiaomi/OPPO/vivo di kamera "Pro Mode" atau "More" bisa menemukan saklar RAW; Huawei di "Pro Mode" aktifkan RAW; Samsung Galaxy S series di mode "Pro" ada opsi RAW. Ukuran file biasanya 20-50MB.
  \item \textbf{Saran Umum}: Foto sehari-hari tidak perlu aktifkan RAW, file terlalu besar memakan ruang; hanya aktifkan di skenario penting (pemandangan, portrait, butuh post-production).
\end{itemize}

\booksubsection{Prinsip Computational Photography}

Fotografi ponsel modern sudah tidak bisa lepas dari “computational photography” — yaitu mengandalkan algoritma bukan hanya lensa dan sensor untuk menghasilkan foto. Memahami beberapa prinsip inti ini membuat Anda lebih paham apa yang dilakukan ponsel saat menekan shutter.

\subsubsection*{Multi-frame fusion:}

Saat ponsel menekan shutter sekali, sebenarnya mengambil banyak foto. Algoritma menyelaraskan multi-frame, fusion berbobot, menghasilkan satu foto berkualitas tinggi.

\begin{itemize}
  \item \textbf{Night Mode}: Dalam 1-3 detik mengambil banyak foto underexpose, setelah diselaraskan ditumpuk — noise saling menetralkan, kecerahan diakumulasi. Hasilnya dalam cahaya redup juga bisa menghasilkan gambar bersih.
  \item \textbf{HDR (High Dynamic Range)}: Memotret underexpose, normal, overexpose tiga (atau lebih) foto, difusion menjadi satu foto yang detail di highlight dan shadow. Skenario backlight sangat efektif.
  \item \textbf{Deep Fusion (Apple) / Photonic Engine}: Untuk skenario cahaya sedang-rendah, analisis multi-frame per piksel lalu difusion, fokus meningkatkan tekstur dan detail, detail seperti rajutan dan rambut akan jauh lebih jelas.
\end{itemize}

\subsubsection*{Semantic segmentation:}

AI tidak hanya memproses seluruh gambar secara seragam, tapi pertama “memahami” apa di dalam gambar — langit, wajah, rumput, bangunan, pakaian — lalu menerapkan strategi pemrosesan berbeda untuk setiap area.

\begin{itemize}
  \item \textbf{Langit}: Secara terpisah meningkatkan biru, menurunkan kecerahan, menghindari overexpose.
  \item \textbf{Wajah}: Identifikasi area warna kulit, menerangkan dan menghaluskan warna kulit secara terpisah, tanpa mengubah detail fitur wajah dan rambut.
  \item \textbf{Rumput/Daun}: Meningkatkan saturasi hijau, meningkatkan sharpening detail.
  \item \textbf{Background Blur}: Mode portrait, AI identifikasi kontur subjek, background blur, subjek jelas.
\end{itemize}

\subsubsection*{Aplikasi deep learning:}

\begin{itemize}
  \item \textbf{AI Scene Recognition}: Kamera otomatis mengidentifikasi skenario "makanan", "hewan peliharaan", "malam", "dokumen", dll. dan mengoptimalkan parameter sesuai. Misalnya mode makanan akan meningkatkan tone hangat dan saturasi, membuat makanan terlihat lebih menggoda.
  \item \textbf{AI Portrait Blur}: Lensa ponsel secara fisik tidak bisa mencapai shallow depth of field seperti DSLR, mengandalkan model deep learning memprediksi tepi subjek, lalu blur background. Algoritma terbaru sudah bisa menangani detail seperti helai rambut dan tepi kaca.
  \item \textbf{AI Erase / AI Expand}: Pilih orang yang lewat atau benda asing di gambar, AI berdasarkan konteks "melengkapi" bagian yang tertutup. Ini aplikasi generative AI di bidang imaging.
\end{itemize}

\subsubsection*{Computational photography vs optical photography:}

\begin{itemize}
  \item \textbf{Keunggulan computational photography}: Volume kecil juga bisa dapat kualitas bagus, skenario ekstrem seperti cahaya redup dan backlight juga bisa menghasilkan foto, pengguna biasa tidak perlu paham fotografi untuk menghasilkan foto "bagus".
  \item \textbf{Kelemahan computational photography}: Gambar kadang "over-processed" — over-sharpening, distorsi warna, warna langit tidak natural; multi-frame fusion saat subjek bergerak akan meninggalkan ghosting; detail "dihitung" bukan direkam nyata.
  \item \textbf{Keunggulan optical photography}: Sensor besar + lensa aperture besar bisa menangkap cahaya secara nyata, kualitas gambar natural, ruang post-production besar; tidak akan menghasilkan artefak karena algoritma "menebak".
  \item \textbf{Tren}: Ponsel flagship modern adalah kombinasi "optik + komputasi" — menggunakan sensor besar untuk mendapatkan informasi nyata, algoritma untuk mengoptimalkan presentasi. Keduanya tidak saling bertentangan.
\end{itemize}

\booksubsection{Jenis Lensa dan Focal Length}

Ponsel flagship modern biasanya dilengkapi tiga sampai empat lensa, masing-masing untuk focal length berbeda. Memahami focal length membuat Anda bisa memilih lensa yang tepat di skenario berbeda.

\subsubsection*{Lensa umum dan focal length ekuivalen:}

\begin{itemize}
  \item \textbf{Ultra wide (13-24mm ekuivalen)}: Bidang pandang terluas, cocok untuk pemandangan, bangunan, ruangan. Bisa menangkap adegan besar yang terlihat, tapi tepi akan distorsi (garis lurus membengkok), saat foto orang wajah akan tertarik. Cocok untuk "adegan besar", tidak cocok untuk "close-up".
  \item \textbf{Kamera utama (24-28mm ekuivalen)}: Lensa yang paling sering dipakai sehari-hari, kualitas terbaik, aperture terbesar. Menyeimbangkan bidang pandang dan kualitas, bisa untuk foto orang dan pemandangan. "1x" di kebanyakan ponsel adalah kamera utama.
  \item \textbf{Telephoto (70-200mm ekuivalen)}: Untuk memotret pemandangan jauh, close-up portrait. Telephoto bisa mengompresi ruang (gunung jauh dan orang dekat terlihat lebih dekat), foto portrait bisa dapat proporsi fitur wajah lebih natural dan background blur.
  \item \textbf{Periscope telephoto (200mm+ ekuivalen)}: Melalui desain jalur cahaya reflektif, dalam ketebalan ponsel mencapai focal length lebih panjang. Bisa memotret detail jelas dari jarak sangat jauh, cocok untuk konser, olahraga, satwa liar. Hanya di model flagship.
\end{itemize}

\subsubsection*{Trade-off ukuran sensor dan ukuran piksel:}

Sensor adalah “fondasi”, menentukan berapa banyak cahaya yang bisa diterima. Industri bilang “sensor lebih besar satu tingkat menekan yang kecil” — dengan jumlah piksel yang sama, sensor lebih besar berarti piksel individual juga lebih besar, cahaya masuk lebih banyak, kualitas gambar lebih baik.

\begin{itemize}
  \item \textbf{Sensor Besar}: Cahaya masuk besar, kualitas foto malam bagus, dynamic range lebar. Kelemahannya ponsel jadi tebal (modul lensa menonjol), biaya tinggi.
  \item \textbf{Piksel Tinggi Sensor Kecil}: Saat siang cahaya cukup, piksel tinggi bisa hasilkan foto lebih jelas, mudah di-crop di post-production. Tapi di cahaya redup setiap piksel kurang cahaya, noise banyak.
  \item \textbf{Pilihan Flagship Mainstream}: Kamera utama kebanyakan 1/1.3 inci sampai 1 inci (sebagian model sudah pakai sensor 1 inci), jumlah piksel sekitar 50 juta — ini titik keseimbangan kualitas dan volume.
\end{itemize}

\subsubsection*{Teknologi pixel binning:}

Untuk menyeimbangkan piksel tinggi siang dan piksel besar malam, produsen ponsel memakai “pixel binning” — menggabungkan beberapa piksel kecil menjadi satu piksel besar untuk output.

\begin{itemize}
  \item \textbf{Quad Bayer (4-in-1)}: Sensor 50 juta piksel, output foto 12.5 juta piksel. Setiap piksel output disintesis dari 4 piksel tetangga, cahaya masuk naik 4x, kualitas gambar lebih bersih.
  \item \textbf{Nonacell (9-in-1)}: Untuk sensor di atas 100 juta piksel, output sekitar 11-12 juta piksel. Prinsip sama, cahaya masuk lebih besar.
  \item \textbf{Efek aktual}: Mode foto default, yang Anda dapatkan adalah foto piksel rendah hasil fusion — kualitas bagus, file kecil. Saat butuh detail ekstrem (seperti foto dokumen, pemandangan jauh), bisa manual beralih ke "High-Res Mode" untuk output resolusi penuh.
\end{itemize}

\booksubsection{Mode Syuting Khusus}

Selain foto sehari-hari, di ponsel ada beberapa mode khusus, jika digunakan dengan benar bisa menghasilkan gambar yang biasanya tidak bisa diambil.

\subsubsection*{Mode panorama:}

\begin{itemize}
  \item Saat memotret ikuti panah petunjuk, gerakkan ponsel secara horizontal perlahan dan stabil, jangan cepat — jika cepat akan muncul misalignment stitching.
  \item Jaga ponsel stabil secara vertikal, hanya geser horizontal. Jika goyang atas-bawah, garis horizon akan bergelombang.
  \item Hati-hati saat ada orang atau mobil bergerak — saat stitching akan muncul "setengah badan" atau "hantu".
  \item Mode panorama juga bisa difoto vertikal — foto vertikal lalu stitch horizontal, cocok untuk memotret bangunan tinggi atau air terjun.
\end{itemize}

\subsubsection*{Slow motion dan time-lapse:}

\begin{itemize}
  \item \textbf{Slow Motion}: Syuting dengan frame rate tinggi 120fps atau 240fps, saat playback menjadi slow motion. Cocok untuk memotret tetesan air, momen olahraga, lompatan hewan. Cahaya harus cukup, mode slow motion waktu masuk cahaya singkat, kualitas cahaya redup buruk.
  \item \textbf{Time-lapse}: Ambil beberapa foto dengan interval lalu disintesis menjadi video fast-forward. Cocok untuk memotret arus lalu lintas, pergerakan awan, sunset, bunga mekar. Durasi syuting dari beberapa menit hingga beberapa jam, harus pakai tripod memfix ponsel, tangan pasti blur.
  \item \textbf{Tip}: Saat time-lapse awan, posisikan ponsel, syuting 15-30 menit, akhirnya dapat video time-lapse sekitar 10 detik, efeknya memukau.
\end{itemize}

\subsubsection*{Mode bintang/astronomi:}

\begin{itemize}
  \item Perlu menjauhi polusi cahaya kota, cari pinggiran kota atau pegunungan tanpa gangguan lampu jalan. Di kota tidak bisa memotret Bima Sakti.
  \item Harus pakai tripod, tangan pasti blur. Ponsel harus stabil ditahan beberapa detik hingga puluhan detik.
  \item Cuaca harus cerah tanpa awan, bulan jangan terlalu besar (bulan purnama akan menutupi cahaya bintang). Malam gelap awal atau akhir bulan kalender lunar paling cocok.
  \item "Night Mode" iPhone akan otomatis mengidentifikasi bintang dan memperpanjang waktu eksposur (maksimal 30 detik); "Star Mode" atau "Astro Mode" flagship Android efeknya serupa. Harus sabar menunggu, setelah menekan shutter jangan gerakkan ponsel.
\end{itemize}

\subsubsection*{Mode macro:}

\begin{itemize}
  \item Untuk memotret bunga, serangga, perhiasan, close-up teks dan benda kecil. Bisa fokus dari jarak 1-2 cm.
  \item Sebagian ponsel punya lensa macro khusus (kebanyakan lensa sensor kecil 2MP atau 5MP), kualitas biasa; flagship memakai lensa ultra wide untuk macro, kualitas lebih baik.
  \item Saat memotret sangat dekat lensa sangat dekat dengan subjek, akan menutupi cahaya, perlu gunakan lampu meja atau senter untuk fill light.
  \item Jarak fokus sangat dekat, depth of field sangat dangkal, subjek maju-mundur beberapa milimeter bisa out of focus. Saat memotret tangan harus stabil, sebaiknya subjek tetap diam.
\end{itemize}

\begin{tcbinfo}{Saran praktis fotografi statis}
\begin{itemize}
  \item Foto harian: HEIF/JPEG default sudah cukup, file kecil, kualitas memadai.
  \item Skenario penting (pemandangan, portrait, butuh post-production): aktifkan RAW, simpan ruang post-production terbesar.
  \item Pilih lensa yang tepat: pemandangan pakai ultra wide, portrait pakai telephoto, harian pakai kamera utama.
  \item Malam dan bintang: harus fix ponsel, aktifkan night mode atau astro mode.
  \item Ponsel sensor besar (seperti model sensor 1 inci) batas atas kualitas lebih tinggi, tapi kebiasaan operasi tetap kunci.
\end{itemize}

\end{tcbinfo}

\begin{tcbpractice}{Coba sendiri: Pahami kemampuan kamera Anda}
\begin{enumerate}
  \item Cek apakah ponsel Anda mendukung syuting RAW (Settings Camera Format), coba aktifkan dan ambil satu foto, bandingkan ukuran file JPEG dan RAW skenario yang sama.
  \item Beralih ke ultra wide, kamera utama, telephoto tiga focal length, masing-masing memotret pemandangan jauh yang sama, rasakan perbedaan sudut pandang dan kualitas.
  \item Cari skenario dengan night mode atau star mode, fix ponsel lalu ambil satu foto malam.
  \item Coba pakai mode panorama foto horizontal adegan besar, perhatikan menjaga ponsel bergerak stabil.
\end{enumerate}

\end{tcbpractice}
